\paragraph{Contact terms in the logarithmic coefficient at \(N_s=6\), \(n_D=2\) (conditional; not closure).}
Keep the setting of the preceding paragraph. To first order in \(1-n\) the surface part of the \(s^0\) heat coefficient of a Laplace-type sector \(-(\nabla^2+E)\) is the cross term of the bulk coefficient \(a_2\) with the conical deltas \(R\to4\pi(1-n)\delta_\Sigma\), \(R_{\mu\nu}\to2\pi(1-n)n_{\mu\nu}\delta_\Sigma\), \(R_{\mu\nu\rho\sigma}\to2\pi(1-n)(n_{\mu\rho}n_{\nu\sigma}-n_{\mu\sigma}n_{\nu\rho})\delta_\Sigma\) (Equations~(55)--(59) of~\cite{Solodukhin2011}) and with any delta \(E_\delta\), \(\Omega_\delta\) of the endomorphism and of the bundle curvature. Without such deltas it reproduces Equation~(70) of~\cite{Solodukhin2011}. For the Hodge operator of a Maxwell field in Feynman gauge with its ghosts, with \(E_\delta\) the Ricci delta and \(\Omega_\delta=2\pi(1-n)\mathcal J\delta_\Sigma\) the Riemann delta (\(\mathcal J\) the rotation of the normal plane), it gives \(S_{\log}=\pi^{-1}[-\tfrac{31}{360}R_\Sigma+\tfrac1{20}W]\), that is the ratios \(62\) and \(12\) of the conformal scalar, and \(-\tfrac{62}{90}\) on a great \(S^2\) of \(S^4\); the contact is \(-(R_\Sigma-R_{ii})/24\pi\), of which \(-R_\Sigma/24\pi\) is the edge-mode density of~\cite{DonnellyWall2016}. The Dirac weights \(-\tfrac12\) on \(\operatorname{tr}(\tfrac R6-E)\), \(1\) on the \(R_{ii}-2R_{ijij}\) term and the spin term of the preceding paragraph reproduce the Dirac \(a\) and \(c\) and the Weyl-covariant form, so \(a=11/360\) is an output of those weights and not a separate input.

For the rotation-invariant self-dual backgrounds with \(\Sigma=\{x^1=x^2=0\}\) and \(g=|\nabla_\Sigma\ln\rho|^2_{g_\Sigma}\), \(g_\Sigma=\rho\,\delta\), one finds at 18 points of random polynomial data (exact rational arithmetic; not proved) \(X'=\tfrac32g\), \(\tau_\Omega=R_\Sigma+\tfrac34g\), so that \(24\pi\sigma_{\log}=\tfrac32g-R_\Sigma\). The statement that the coefficient is an Euler density holds only where \(g=2R_\Sigma\), as on the example above. With all sectors glued periodically and no deltas the coefficient vanishes for every Poisson bivector at \((6,2)\), because \(N_{\mathrm{tot}}=0\) and the weighted \(a_1\) vanishes; the value \(\sigma_{\log}\) is the difference between the antiperiodic and the periodic Dirac gluing.

\emph{Contacts.} On the family \(\Theta^{12}=1/f(x^1,x^2)\), \(\Theta^{34}=1\), with \(\ln f=(n-1)\ln|z|\), the scalar, connection and Dirac data contain \(\tfrac12\Box_0\ln f\) with the same coefficient in every component, so that \(E_\delta=-\pi(1-n)\delta_\Sigma\) uniformly (the contact with \(Q=\tfrac14\) above), and the connection bundle curvature contains \(\Omega_{12,\delta}=2\pi(1-n)\mathcal J\delta_\Sigma\) (symbolic computation). Let \(C_b\) be the logarithmic cross term of \(E_\delta\) on the six scalars, the four components of \(a^\mu\) and the ghosts, and \(C_\Omega=\operatorname{tr}(\mathcal J\Omega^{v}_{12})/24\pi\) that of the connection bundle. Then, in units \(1/24\pi\), \(24\pi\sigma_{\log}+24\pi C_b=-\tau_\Omega\) and \(\operatorname{tr}_V(\mathcal J\Omega^v_{12})=\tau_\Omega\) (the vector is \((\tfrac12,\tfrac12)\) and the Dirac field \((\tfrac12,0)\oplus(0,\tfrac12)\) under \(\mathfrak{su}(2)\oplus\mathfrak{su}(2)\)), hence \(\sigma_{\log}+C_b+C_\Omega=0\). This was checked at random closed two-forms (not self-dual, \(\mathfrak J\neq0\)) and on the \(\mathfrak J=0\) families above. On the example, \(C_b+C_\Omega=-R_\Sigma/12\pi\), which cancels \(\sigma_{\log}\); \(C_b\) alone gives \(-3R_\Sigma/16\pi\), \(C_\Omega\) alone \(5R_\Sigma/48\pi\). A Herzog--Nishioka term removing \(E_\delta\)~\cite{HerzogNishioka2016} leaves \(\sigma_{\log}+C_\Omega=3g/32\pi\ne0\). An edge density of the Donnelly--Wall type for a field with coupling \(w=\rho^{-1}\) is \(-(4\pi)^{-1}(R_\Sigma/6-\tfrac12\Delta\ln w-\tfrac14|\nabla\ln w|^2)\), and \(\sigma_{\log}\) plus it does not vanish. Dirac fields and ghosts carry no edge mode and no further contact; the \(Q\) contacts of the scalars and ghosts vanish at logarithmic order when \(R[g_\rho]=0\), and with them the \(\LUV^2\) coefficient is \(c_A=0=c_T\).

These statements are for the Poisson-limit operators, one loop, \(\alpha=1\), \(N_s=6\), \(n_D=2\), a totally geodesic \(\Sigma\), cross terms at first order in \(1-n\) (the \(n<1\) branch of the rotated connection index of the paragraph on replica prescriptions), and a replica whose connection index is transported with the replica angle and whose bosonic endomorphisms carry the delta above; the non-delta singular first-order parts of the smooth family, the \(n>1\) branch and the state selected by the matrix model are not treated. The fixed-state replica without these contacts keeps the difference of the preceding paragraph. They give no value or sign of \(G_{\mathrm{eff}}\). Problem~B remains open and PROGRAMME-CONDITIONAL.
